\section{Related work}
\label{sec:related}

\paragraph{Runtime-defined label spaces.} Deciding over a candidate set supplied at inference time is an
old problem. GILE embeds inputs and label descriptions in a shared space so unseen labels are scorable
\citep{pappas2019gile}; entailment reformulation turns arbitrary label names into hypotheses
\citep{yin2019zeroshot}; TARS conditions a sentence encoder on the label text itself
\citep{halder2020tars}; GLiNER and GLiNER2 encode text and a runtime schema jointly in one bidirectional
pass \citep{zaratiana2023gliner,zaratiana2025gliner2}, the latter also appearing on the JevBench
leaderboard we report against. The cost axis is the classical cross- versus dual-encoder trade-off: a
cross-encoder reranker attends candidate to input and pays one forward pass per candidate
\citep{nogueira2019bertrerank}, a dual encoder is candidate-independent but interacts only at a dot
product \citep{reimers2019sbert}, and late-interaction and decomposition methods sit between --
Poly-encoders \citep{humeau2020polyencoders}, ColBERT \citep{khattab2020colbert}, DeFormer
\citep{cao2020deformer}, PreTTR \citep{macavaney2020prettr}, LUMEN \citep{dejong2023lumen} -- with
compiling a cross-encoder into a cheaper student as standard practice \citep{lu2022erniesearch}.
Decoder LLMs with a classification head close the loop \citep{li2023labelsupervised}. Four distinctions
are worth making explicit. Against \emph{fixed classifiers and sequence classification}: our candidates
are runtime-defined text, not fixed class indices. Against \emph{GILE, TARS, and dual encoders}: our
central negative result is precisely that a candidate-independent representation is insufficient for the
targeted question-conditioned closed-book regime (\Cref{sec:repr:blind}). Against \emph{cross-encoders
and rerankers}: Typical reuses one state key-value prefix across many decisions and emits a typed
distribution with an abstention event, rather than independently reranking state--candidate pairs.
Against \emph{poly-encoders, ColBERT, and late interaction}: our candidates are rendered inside the
causal computation and acquire contextual slot and set information that a late-interaction embedding of
the candidate in isolation does not carry. Our own energy-scoring shortlisting front end
(\Cref{sec:train:latency}) is a member of this family and is retained for exactly the sub-task --
bounded-$\numcands$ pre-filtering -- that does not require question-conditioned capability.

\paragraph{Reading the decision out of hidden states instead of generating it.} This is our closest prior
art. Linear probes and their controls establish that hidden states carry decodable structure
\citep{alain2017probes,hewitt2019controltasks,belinkov2022probing}; layerwise analysis shows different
depths carry different abstractions \citep{tenney2019pipeline,belrose2023tunedlens}; states anticipate
tokens not yet emitted \citep{pal2023futurelens,wu2024planahead}; and MCQ-specific circuit analysis
localizes the symbol-binding step \citep{lieberum2023circuit}. Most directly, \citet{wong2026models} show
a two-stage computation in ordinary multiple-choice prompting: the winning option is decodable in content
space at the option boundary \emph{before} the answer symbol is bound, a finding traced across languages
by \citet{bhatt2026representation}. The letter interface itself is documented as unreliable: selection
bias and option-order sensitivity \citep{zheng2024selectionbias,pezeshkpour2024optionorder}, symbol-identity
effects and binding fixes \citep{yang2025option,xue2024symbolbinding}, first-token probabilities that
disagree with the generated answer \citep{wang2024firsttoken}, and the choices-only control quantifying
how much multiple-choice accuracy needs no question at all \citep{balepur2024artifacts} -- the control we
adopt as our primary diagnostic (\Cref{sec:experiments:metrics}). Prior work observes the latent decision
and diagnoses the symbol interface. We do not claim novelty at the level of ``the answer is latent before
generation'' -- that is \citet{wong2026models}. What we contribute is to make the candidate-aware latent
representation the \emph{deployed output interface} and to measure what is lost when the candidates are
removed from the computation, including that the last layer is the wrong tap
(\citealp{gromov2024deeperlayers,men2024shortgpt,schwartz2020righttool}; \Cref{sec:repr:depth}) and that
the readout must be over contextual option spans, not slot-agnostic embeddings
(\Cref{sec:repr:readout}).

\paragraph{Calibration, distillation, and what survives a change of interface.} Temperature scaling
largely fixes in-distribution expected calibration error \citep{guo2017calibration}, though binned ECE is
itself a biased estimator \citep{naeini2015calibration,kumar2019verifiedcalibration}, motivating proper
scoring rules as the honest target \citep{gneiting2007scoringrules,brier1950verification}, with the ranked
probability score as their ordinal-aware member \citep{epstein1969rps,murphy1970rps}. Training-time levers
include label smoothing \citep{szegedy2016labelsmoothing,mueller2019labelsmoothing} and soft-target
distillation, which transfers more than argmax \citep{hinton2015distillation,kim2016seqkd} and improves
probability quality even without a stronger teacher under the label-smoothing view of KD
\citep{yuan2020teacherfree,zhang2020selfdistillation}; fine-tuning specifically can distort calibration
through its interaction with prior knowledge \citep{wang2025cogcalib,desai2020calibration}. Our
calibration claims are training-time rather than post-hoc: soft and ordinal-smoothed targets
(\Cref{sec:method:typed}), a Brier term in the objective, and checkpoint selection on a held-out
calibration metric (\Cref{sec:method:selection}) rather than a single global temperature fit. Our
frozen-teacher distillation term (\Cref{sec:boundary:frozen}) is a distillation target in this sense but a
narrow one -- the teacher is the same backbone read with a fixed few-shot rendering, so the relevant
question is componential: which parts of a teacher's behavior survive a change of interface. Our answer
(priors and calibration, not question-conditioned knowledge) is the negative half of this literature,
consistent with context distillation compiling a prompt into weights
\citep{askell2021assistant,snell2022contextdistillation}.

\paragraph{Prefix sharing and the cost of removing the decode loop.} Sharing a long prefix across many
suffixes is established systems work \citep{juravsky2024hydragen,zheng2023sglang,kwon2023pagedattention,
gim2023promptcache,pope2023scaling,dao2022flashattention,cheng2023batch}; generation-side accelerators are
an informative contrast, since speculative decoding and Medusa exist because each sequential decode step
is overhead-bound rather than compute-bound
\citep{leviathan2023speculative,chen2023speculativesampling,cai2024medusa}. We contribute no new kernel; we
report the measured, single-stream-harness consequence of deleting the decode loop entirely
(\Cref{sec:train:latency}).

\paragraph{Is generation necessary? Expressivity and latent reasoning.} Chain-of-thought buys serial
computation that a constant-depth transformer provably lacks in one forward pass
\citep{wei2022cot,li2024cotserial,merrill2024expressive,merrill2023parallelism,feng2023chain}, and helps
empirically mainly on math and symbolic tasks \citep{sprague2025cot}; filler and pause tokens show part of
the gain is extra computation rather than verbalized reasoning
\citep{goyal2024pause,pfau2024filler}, and implicit or latent chain-of-thought internalizes the steps
\citep{deng2023implicitcot,deng2024internalize,hao2024coconut}. This literature delimits our scope
(\Cref{sec:discussion}): a single non-generative forward pass is the right instrument for bounded,
evidence- or prior-driven decisions, and is predicted to fail on serial composition -- exactly where our
hard tier (temporal, probability, trade-off families) remains near chance.

\paragraph{Ordinal regression, abstention, and long-context failure modes.} Our Score primitive is the
minimal member of the ordinal-regression family
\citep{mccullagh1980ordinal,frank2001ordinal,niu2016ordinalcnn,cao2020coral,shi2023corn,diaz2019softordinal}:
a soft, ordinal-smoothed target that changes no decision and fixes probability quality, where a
cumulative-link head did not win at matched budget (\Cref{sec:train:typed}). Our abstain option is
architectural rather than a classical reject rule
\citep{chow1970reject,geifman2017selective,geifman2019selectivenet,hendrycks2017msp,liu2020energyood}, with
out-of-scope intent detection as the applied benchmark \citep{larson2019clinc,casanueva2020banking77}. On
long context, retrieval position and length are known to change and degrade answers
\citep{liu2024lostmiddle,levy2024moretokens,shi2023distracted,hsieh2024ruler}; our truncation finding
(\Cref{sec:boundary:truncation}) is the training-side analogue, an instance of data-pipeline failures
propagating silently through ML systems \citep{sambasivan2021datacascades}.

\paragraph{Fine-tuning versus frozen backbones with in-context examples.} Fine-tuning can distort
pretrained features and underperform out of distribution \citep{kumar2022finetunedistort}, fair
comparisons of few-shot fine-tuning against in-context learning are close \citep{mosbach2023fewshot}, LoRA
both learns and forgets less than full fine-tuning \citep{biderman2024loralearnsless,hu2022lora} yet
continual fine-tuning still forgets \citep{luo2025forgetting}, and compositional-generalization benchmarks
show where learned grammars stop composing out of distribution
\citep{lake2018scan,keysers2020cfq,kim2020cogs,hupkes2020compositionality}. Two of our findings live here:
rubric-conditioned training transfers inside its own rule grammar and not to unseen composition
(\Cref{sec:boundary:hard}), and our trained 14B leads its frozen self on the standard tier while a
same-size frozen backbone with three exemplars still leads on the hard tier
(\Cref{sec:boundary:frozen}).

\paragraph{Typed decision benchmarks and commercial systems.} JevBench, one of the evaluation suites we
use, targets a class of commercial ``System-One'' decision models -- most visibly TypeSafe's Jev,
documented only by vendor pages and a third-party black-box probe, closed-weight and closed-data. We
therefore say that Typical targets a similar interface and task class, not that it reproduces Jev's
architecture, which we cannot inspect. The leaderboard carries several open replications that publish
weights and scores but, apart from GLiNER2 \citep{zaratiana2025gliner2} and the ModernBERT backbone used
by Laya \citep{warner2024modernbert}, no method papers.

\paragraph{What we claim as new.} Not ``reading decisions from hidden states'', which has ample prior art
above. The defensible contribution is narrower and sharper: we study the design space for turning a
pretrained causal LM into a reusable typed decision primitive under a cached-state constraint, and
identify a specific empirical boundary in it -- candidate-blind reusable decision states fail to preserve
question-conditioned capability, while candidate-contextual readout succeeds -- together with the
diagnostic ($\dq$) that makes the boundary visible at all. A secondary methodological contribution is
that rendering is a first-class experimental factor for this task class: a frozen Qwen3.5-9B checkpoint
moves 12.5 points of standard-tier accuracy on rendering choice alone (\Cref{sec:boundary:rendering}).
We draw the cautious conclusion from that, not the convenient one: it means cross-entry benchmark
comparisons in this area conflate model capability with interface compatibility -- \emph{ours
included} -- and that rendering should be held fixed when architectures are compared, not that any
particular published result is wrong.
